
\documentclass[UTF8,a4paper,11pt]{ctexart}

% ============================================================
% 页面与基础宏包
% ============================================================
\usepackage[
    top=2.20cm,
    bottom=2.20cm,
    left=2.30cm,
    right=2.30cm,
    headheight=15pt
]{geometry}

\usepackage{amsmath}
\usepackage{booktabs}
\usepackage{longtable}
\usepackage{tabularx}
\usepackage{array}
\usepackage{enumitem}
\usepackage{fancyhdr}
\usepackage{lastpage}
\usepackage{hyperref}
\usepackage{xcolor}
\usepackage[most]{tcolorbox}

\hypersetup{hidelinks}

% ============================================================
% 黑白正式配色
% ============================================================
\definecolor{BoxGray}{HTML}{F5F5F5}
\definecolor{LineGray}{HTML}{B8B8B8}
\definecolor{TextGray}{HTML}{555555}

% ============================================================
% 正文格式
% ============================================================
\setlength{\parindent}{2em}
\setlength{\parskip}{0pt}
\linespread{1.18}

\setlist[itemize]{
    leftmargin=2em,
    itemsep=0.10em,
    topsep=0.20em,
    parsep=0pt
}

\setlist[enumerate]{
    leftmargin=2.2em,
    itemsep=0.10em,
    topsep=0.20em,
    parsep=0pt
}

\renewcommand{\arraystretch}{1.32}

% ============================================================
% 标题格式
% ============================================================
\ctexset{
    section = {
        format = \large\bfseries,
        beforeskip = 1.30em,
        afterskip = 0.55em
    },
    subsection = {
        format = \normalsize\bfseries,
        beforeskip = 0.95em,
        afterskip = 0.35em
    }
}

% ============================================================
% 页眉页脚
% ============================================================
\pagestyle{fancy}
\fancyhf{}

\fancyhead[L]{\small AI工具使用详情}
\fancyhead[R]{\small 2026年全国大学生数学建模竞赛}
\fancyfoot[C]{\small \thepage/\pageref{LastPage}}

\renewcommand{\headrulewidth}{0.4pt}
\renewcommand{\footrulewidth}{0pt}

% ============================================================
% 表格列
% ============================================================
\newcolumntype{Y}{>{\raggedright\arraybackslash}X}
\newcolumntype{L}[1]{>{\raggedright\arraybackslash}p{#1}}
\newcolumntype{C}[1]{>{\centering\arraybackslash}p{#1}}

% ============================================================
% 正式交互记录框
% ============================================================
\newtcolorbox{recordbox}[2]{
    enhanced,
    breakable,
    colback=white,
    colframe=LineGray,
    boxrule=0.65pt,
    arc=0pt,
    outer arc=0pt,
    left=4mm,
    right=4mm,
    top=3.2mm,
    bottom=3.2mm,
    before skip=0.7em,
    after skip=0.8em,
    title={
        \textbf{#1}
        \hfill
        \normalfont\small #2
    },
    coltitle=black,
    colbacktitle=BoxGray,
    fonttitle=\normalsize,
    boxed title style={
        boxrule=0pt,
        arc=0pt
    }
}

\newcommand{\recorditem}[1]{%
    \par\vspace{0.28em}
    \noindent\textbf{#1}\quad
}

% ============================================================
% 正文开始
% ============================================================
\begin{document}

% ============================================================
% 标题
% ============================================================
\begin{center}
    \vspace*{-0.7cm}

    {\zihao{2}\bfseries AI工具使用详情}

    \vspace{0.30cm}

    {\large 2026年全国大学生数学建模竞赛}

    \vspace{0.45cm}

\end{center}

\vspace{0.45cm}

本项目围绕含光伏、负荷、储能和外购电的微网调度问题展开。
建模过程中使用的生成式人工智能工具均来自 DeepSeek，
具体使用了
\textbf{DeepSeek V4.1 Flash}
和
\textbf{DeepSeek V4 Pro}
两个模型版本。

其中，DeepSeek V4 Pro 主要用于需要较强综合推理的
备选方案比较、参数分析、信息可用范围检查和策略评价；
DeepSeek V4.1 Flash 主要用于代码实现检查、
局部逻辑排查以及论文文字审阅等高频辅助任务。

团队负责数学模型建立、核心公式推导、程序实现、参数最终确定、
实验运行、模型比较与最终选择以及最终结论形成。

团队使用AI时，一般先独立形成模型、程序或实验结果，
再根据任务性质选择相应模型，
将具体问题、必要的代码片段、公式或实验现象提交给AI辅助检查。
AI输出均作为参考意见，
最终是否采用由团队根据题目条件、数学推导、程序运行和实验结果判断。

% ============================================================
\section{AI工具使用概况}

本项目混合使用 DeepSeek V4.1 Flash 和 DeepSeek V4 Pro。
其中，需要综合比较、较长推理和建模判断的任务主要使用 V4 Pro；
代码检查、局部问题排查和文字审阅等任务主要使用 V4.1 Flash。

\begin{center}
\small
\begin{tabularx}{\textwidth}{
    L{0.22\textwidth}
    L{0.24\textwidth}
    Y
}
\toprule
\textbf{模型版本} &
\textbf{主要使用阶段} &
\textbf{主要用途} \\
\midrule

DeepSeek V4 Pro &
建模方案比较、参数分析及策略评价 &
比较简单预测方法与复杂预测模型；
分析风险参数的统计含义；
检查预测数据和日内预报的可用范围；
讨论价格预测指标与最终调度费用之间的关系；
对关键建模选择提供备选分析思路。 \\

DeepSeek V4.1 Flash &
程序调试、快速检查及论文整理 &
检查 Python 程序中的时间索引、SOC递推、
滚动窗口和边界处理；
辅助排查局部代码逻辑；
审查论文中的结果表述、术语使用和因果措辞。 \\

\bottomrule
\end{tabularx}
\end{center}

两个模型均未直接生成最终模型或决定最终方案。
对于涉及模型选择、参数确定和程序修改的关键内容，
团队均重新通过题目原文、数学推导、程序实验或手工计算进行核验。
以下选取较有代表性的使用过程进行说明。

% ============================================================
\section{典型AI使用与交互记录}

% ============================================================
\begin{recordbox}
{交互记录 01：预测模型选择与复杂模型比较}
{问题二 · DeepSeek V4 Pro}

\recorditem{使用目的}
团队已经构造 F0、F1、F2 三种只使用目标日前历史数据的预测方法，
同时考虑是否进一步使用 LSTM、Transformer 等复杂时序模型。
需要判断复杂模型是否确有必要，
并避免因提前使用未来数据或不同预测方法的比较条件不一致而产生不公平比较。

\recorditem{主要提示内容}
团队向 DeepSeek V4 Pro 说明数据按日、按10分钟记录，
并明确预测时只能使用目标日之前已经获得的信息。
要求其比较简单预测方法和复杂时序模型时，
重点分析未来信息使用风险、过拟合风险、比较条件的一致性，
以及预测精度与最终调度费用之间的关系，
而不直接替团队确定最终模型。

\recorditem{AI反馈摘要}
DeepSeek V4 Pro 建议保证所有预测方法使用相同的历史信息，
并在相同数据区间内比较预测误差、应急购电量和最终购电费用。
同时指出，预测误差更低不一定意味着整个调度系统费用更低；
复杂模型只有在相同条件下能够稳定改善最终调度结果时，
其额外复杂度才具有实际价值。

\recorditem{团队处理}
团队采纳了“使用相同历史信息、在相同数据区间内比较，
并结合最终调度结果评价”的原则，
但没有直接采用AI对具体预测模型的选择意见。

\recorditem{人工核验}
团队重新逐日检查预测数据的截止时间，
并在统一历史信息和比较区间下运行不同预测方案。
比较预测误差、应急购电量和全年购电费用后发现，
复杂模型没有表现出稳定的最终费用优势，
因此最终没有采用 LSTM 或 Transformer 作为主预测模型。

\end{recordbox}

% ============================================================
\begin{recordbox}
{交互记录 02：残差分位参数与敏感性分析}
{问题二 · DeepSeek V4 Pro}

\recorditem{使用目的}
问题二利用历史净负荷预测误差的经验分位数确定风险裕量。
团队需要确定历史窗口 $W$ 和分位参数 $\alpha$，
并正确理解 $\alpha$ 的统计含义。

\recorditem{主要提示内容}
团队向 DeepSeek V4 Pro 询问
$\alpha=0.8$、$0.9$ 和 $0.95$
分别表示什么，
以及该参数是否能够直接理解为“整日可靠度”。
同时要求AI提出值得进一步检查的敏感性因素。

\recorditem{AI反馈摘要}
DeepSeek V4 Pro 指出，
$\alpha$ 对应历史预测残差分布中的经验分位位置，
不能直接解释为一天内全部时段同时满足要求的联合概率。

AI同时建议进一步检查储能效率、储能容量、
充放电功率、初始SOC、终端状态以及预测误差放大程度
对最终结果的影响。

\recorditem{团队处理}
团队采纳了AI对经验分位含义的解释，
并按照建议扩展敏感性分析；
但 $W$ 与 $\alpha$ 的最终取值没有直接采用AI建议，
而是由全年逐日模拟结果确定。

\recorditem{人工核验}
团队重新比较不同 $W$ 和 $\alpha$ 组合，
并分别改变储能效率、容量、功率、初始SOC和预测误差倍率。
最终依据全年购电费用、应急购电量及结果稳定性确定最终参数组合。

\end{recordbox}

% ============================================================
\begin{recordbox}
{交互记录 03：SOC递推与滚动优化程序排查}
{问题二至四 · DeepSeek V4.1 Flash}

\recorditem{使用目的}
滚动优化程序初版已经能够运行，
但团队检查代表日结果时发现，
部分时段的SOC变化与手工计算结果不一致，
因此需要进一步定位程序实现问题。

\recorditem{主要提示内容}
团队将SOC递推公式和相关 Python 代码片段提交给
DeepSeek V4.1 Flash。

当 $x_t$、$y_t$ 分别表示交流侧充电量和放电量时，
团队使用的SOC递推关系为
\[
S_{t+1}
=
S_t+0.9x_t-\frac{y_t}{0.9}.
\]

同时要求AI重点检查：

\begin{itemize}
    \item 充放电效率方向是否一致；
    \item 时间索引是否错位；
    \item 首时段执行是否正确；
    \item 滚动窗口末端是否存在边界问题；
    \item SOC达到上下限时如何修正充放电量；
    \item 应急购电是否可能被错误地用于主动充电。
\end{itemize}

\recorditem{AI反馈摘要}
DeepSeek V4.1 Flash 提示团队重点检查
滚动窗口边界数组与实际执行时段之间的对应关系，
并指出按照SOC上下限修正实际充放电量后，
可能出现“优化得到的充放电量”与“实际执行后的SOC变化”不一致的问题。

\recorditem{团队处理}
团队根据提示重新检查代码，
确认滚动边界数组和SOC边界处理确有需要修正的部分，
随后重新编写相关边界处理程序。
AI给出的代码没有直接复制使用。

\recorditem{人工核验}
修改后，团队构造3个连续时段进行手工计算，
逐步核对SOC变化，
随后逐时段检查母线能量平衡，
并对全年结果运行
\texttt{balance}、
\texttt{SOC chain}、
\texttt{overlap} 和
\texttt{settlement}
等程序检查。
手工计算与程序结果一致后才保留修改。

\end{recordbox}

% ============================================================
\begin{recordbox}
{交互记录 04：日内预报更新与合同调整时序检查}
{问题三 · DeepSeek V4 Pro}

\recorditem{使用目的}
题目只在
0:00、6:00、12:00 和 18:00
提供新的光伏预测信息。
团队需要确认程序严格遵守四个规定的预报时点，
并避免已经执行的时段被重新优化或重复结算。

\recorditem{主要提示内容}
团队向 DeepSeek V4 Pro 说明四个预测发布时间和初版合同调整逻辑，
要求重点检查：

\begin{itemize}
    \item 程序是否在非规定预报时点读取新的预测信息；
    \item 是否存在使用未来信息的可能；
    \item 已经执行的时段是否保持不变；
    \item 当前实际SOC是否正确传递到下一次优化；
    \item 同一目标时段的合同差额是否可能重复结算。
\end{itemize}

\recorditem{AI反馈摘要}
DeepSeek V4 Pro 指出，
如果程序在每一个10分钟时段都根据最新信息重新生成剩余合同，
实际上相当于连续获得新的预测信息，
与题目只允许四次预报更新的条件不一致。

AI建议只在规定的预报时点更新购电计划：
每次仅在新的光伏预报发布后重新计算尚未执行时段的购电计划；
已经执行的时段不再调整，已发生费用不再改变；
剩余时段根据当前SOC和当前已知信息重新优化。

\recorditem{团队处理}
团队重新核对题目原文后确认这一问题存在，
因此否决连续更新的实现方式，
将程序改为仅在四个规定预报时点调整尚未执行时段的购电计划。

\recorditem{人工核验}
团队分别抽查
0:00、6:00、12:00 和 18:00
四个预报时点前后的数据版本、SOC和购电计划状态；
重新逐时段累计费用，
确认已执行时段不会被重新修改，
并检查每个目标时段只相对原合同进行一次最终差额结算。
此外，通过多组不同预报组合对比实验检查不同预报更新方式下的结果变化。

\end{recordbox}

% ============================================================
\begin{recordbox}
{交互记录 05：波动电价下的评价方式检查}
{问题四 · DeepSeek V4 Pro}

\recorditem{使用目的}
团队比较昨日同一时段、近7日均值、
近30日均值和近30日中位数等价格预测方式时发现，
价格 MAE、RMSE、尖峰误差与最终购电费用的排序并不完全一致，
因此需要确定合理的策略评价方式。

\recorditem{主要提示内容}
团队要求 DeepSeek V4 Pro 解释：
为什么价格预测误差更小，
不一定意味着实际购电费用更低；
同时要求其检查预测价格与实际结算电价在模型中的作用是否应严格区分。

\recorditem{AI反馈摘要}
DeepSeek V4 Pro 指出，
预测模型与决策模型的评价目标并不相同。
价格预测误差只反映预测值与实际电价之间的偏差，
而最终调度费用还受到误差发生时段、
储能状态、合同调整能力和应急购电等因素影响。

因此AI建议除 MAE、RMSE 外，
还同时比较尖峰价格误差、应急购电量和全年购电费用，
并将提前知道全天实际电价的情况仅作为不可实施的参考。

\recorditem{团队处理}
团队采纳多指标评价方式，
没有以价格MAE最小作为唯一的策略选择依据，
同时在程序中严格区分预测电价与实际结算电价。

\recorditem{人工核验}
团队重新检查价格数据时间标签及程序调用位置，
确认预测阶段没有读取未来实际电价；
随后对不同价格预测方法分别运行全年逐日模拟，
综合比较价格误差、尖峰误差、应急购电量和购电费用后确定最终方案。

\end{recordbox}

% ============================================================
\begin{recordbox}
{交互记录 06：结果表述与论文文字审查}
{全文 · DeepSeek V4.1 Flash}

\recorditem{使用目的}
部分实验结果之间存在明显相关关系，
但团队希望避免将特定实验条件下的现象写成过强的普遍因果结论。

\recorditem{主要提示内容}
团队将部分结果分析文字提交给 DeepSeek V4.1 Flash，
要求检查是否存在：
将相关关系直接写成因果关系、
将单组实验结果过度推广，
或者将预测指标与调度效果简单等同的情况。

\recorditem{AI反馈摘要}
DeepSeek V4.1 Flash 建议减少
“必然导致”“显著提高”等过强措辞，
将其改写为带有实验条件的描述，
例如“在当前策略和信息条件下费用下降”
或“该现象与某项变化同时出现”。

\recorditem{团队处理}
团队参考AI意见重新修改相关文字，
但没有直接复制AI生成的完整段落。

\recorditem{人工核验}
修改后的结论逐项与预报组合对比、
敏感性分析和实际实验结果对应，
确保论文表述没有超出实验结果可以支持的范围。

\end{recordbox}

% ============================================================
\section{AI建议的采纳、修改与否决情况}

为明确AI输出对最终方案的实际影响，
主要建议及团队处理情况汇总如下。

\begin{longtable}{
    L{0.31\textwidth}
    C{0.15\textwidth}
    L{0.43\textwidth}
}
\toprule
\textbf{AI建议或输出} &
\textbf{团队处理} &
\textbf{人工核验依据} \\
\midrule
\endfirsthead

\toprule
\textbf{AI建议或输出} &
\textbf{团队处理} &
\textbf{人工核验依据} \\
\midrule
\endhead

保证不同预测模型使用相同历史信息并在相同数据区间内比较 &
采纳 &
逐日检查数据截止时间，并在相同数据区间重新计算预测误差、
应急购电量和购电费用 \\

利用历史残差经验分位数进行风险修正 &
修改后采纳 &
重新比较不同 $W$ 和 $\alpha$，并进行多组敏感性分析 \\

使用 LSTM、Transformer 等复杂模型作为主预测器 &
验证后否决 &
在相同历史信息和全年逐日模拟条件下，
没有获得稳定的最终费用优势 \\

检查滚动优化中的SOC边界和时段边界处理 &
修改后采纳 &
通过3时段手算、逐时段母线平衡及全年程序检查后确认修改有效 \\

连续获取新光伏预报并持续调整合同 &
否决 &
与题目规定的四个预报更新时间不一致，
存在引入未来信息的风险 \\

仅根据价格MAE最小选择问题四策略 &
未采纳 &
价格预测误差与最终购电费用不完全等价，
需要结合应急购电量和最终购电费用评价 \\

使用过强的因果关系表述 &
修改 &
结合预报组合对比和敏感性分析重新核对后，
改为限定实验条件的结果表述 \\

\bottomrule
\end{longtable}

% ============================================================
\section{人工核验与AI参与边界}

AI输出仅作为建模过程中的辅助意见使用。
对于涉及模型选择、参数设置、程序修改和结果解释的内容，
团队均在采用前重新核对题目条件、数学关系和程序运行结果。
人工核验主要围绕
“题目条件是否满足、数学关系是否正确、程序实现是否一致、
实验结论是否稳定”
几个方面展开。

\begin{center}
\small
\begin{tabularx}{\textwidth}{
    L{0.20\textwidth}
    Y
    Y
}
\toprule
\textbf{核验内容} &
\textbf{主要检查方式} &
\textbf{对应目的} \\
\midrule

题目与可用信息 &
重新对照题目原文和附件，
检查预测数据可获得时间、
问题三四个预报发布时间、
合同调整规则以及问题四价格数据的使用范围 &
防止在预测、合同调整或价格决策中提前使用未来信息 \\

数学关系 &
重新推导功率与电量换算、母线平衡、SOC递推、
应急购电和合同差额结算等核心关系 &
确认AI对公式或程序的解释与实际模型一致 \\

程序实现 &
利用小规模手算样例检查关键递推，
并逐时段核对变量索引、时间步长、跨日SOC、
预测窗口、已执行时段状态和费用累计 &
确认程序修改没有引入新的边界错误或重复计算 \\

模型与参数 &
在相同历史信息和相同数据区间下重新运行全年逐日模拟，
比较不同预测模型、风险参数、储能参数及预报更新组合 &
避免依据单次AI建议直接确定模型或参数，
并检查最终结论是否稳定 \\

结果一致性 &
将逐时段CSV、逐日汇总数据、Excel工作簿和论文表格相互核对，
并重点检查问题三、问题四中的合同差额结算 &
保证论文中使用的最终结果与实际程序输出一致 \\

\bottomrule
\end{tabularx}
\end{center}

\vspace{0.5em}

在上述过程中，
DeepSeek V4 Pro 和 DeepSeek V4.1 Flash
分别承担不同类型的辅助工作，
而模型构建、程序实现和最终判断始终由团队完成。
具体参与边界如下。

\begin{center}
\small
\begin{tabularx}{\textwidth}{
    L{0.22\textwidth}
    Y
}
\toprule
\textbf{责任主体} &
\textbf{主要工作内容} \\
\midrule

DeepSeek V4 Pro &
主要用于需要综合推理的任务，
包括备选预测和建模方案比较、
风险参数含义分析、
预测数据与预报更新时点检查、
日内合同调整逻辑分析、
价格预测与调度费用评价关系讨论等。 \\

DeepSeek V4.1 Flash &
主要用于高频辅助检查，
包括 Python 程序中的索引、SOC递推、
滚动窗口和边界逻辑排查，
以及结果文字、术语和因果表述的审阅。 \\

团队 &
建立整体数学模型并推导核心公式与约束；
完成 Python 程序的最终实现；
确定模型和参数；
运行全部数值实验和全年逐日模拟；
比较不同方案并作出最终模型选择；
核实最终结果并形成论文结论。 \\

\bottomrule
\end{tabularx}
\end{center}

AI提出的内容并未全部进入最终方案。
其中，DeepSeek V4 Pro 的输出主要用于帮助团队扩展备选方案和检查建模逻辑，
DeepSeek V4.1 Flash 的输出主要用于辅助发现程序实现和文字表达中的局部问题。
无论使用哪一个模型版本，
AI输出均不直接作为模型选择、参数取值或实验结论的依据。

例如，复杂预测模型、连续更新光伏预报以及仅根据价格MAE选择策略等做法，
均在团队重新检查题目条件和实验结果后被否决；
关于风险参数、SOC边界处理和评价指标等建议，
也均经过重新计算、程序测试或敏感性分析后才被采用。

因此，本项目中AI的作用主要体现在
\textbf{提供备选思路、辅助发现问题和协助检查实现}，
没有替代团队进行核心建模、程序实验和最终决策。
报告中的模型、参数、数值结果和最终结论均由团队自行核实并负责。

\vfill

\begin{center}
\small
\textit{AI输出均经团队人工复核后使用，最终模型、实验结果和论文结论由团队负责。}
\end{center}

\end{document}

